\section{从语言模型到聊天机器人}

\subsection{原始语言模型的行为}

讲师使用 Google 较早期的 LaMDA 模型做演示，因为它还没有经过大量 post-training 微调，更容易观察到原始的 next-word prediction 行为。

当直接输入 ``Hi, do you have any recommendations for dinner?'' 时，模型的输出远不像一个合格的聊天机器人。它确实提到了一些餐厅推荐（如 ``The Fat Duck''），但随后输出了 ``TripAdvisor staff removed this post'' 这样毫无关系的文字。

\begin{importantbox}{语言模型是训练数据的``模糊查询''}
语言模型本质上是在做其训练数据的模糊查询（fuzzy lookup）。LaMDA 生成 ``TripAdvisor staff removed this post'' 很可能是因为训练数据中包含 TripAdvisor 论坛内容，而这句标准化文字在论坛帖子中出现频率极高，因此模型倾向于生成它。这不是模型在``思考''，而是在重现训练数据的统计模式。
\end{importantbox}

\subsection{逐步构建聊天界面}

讲师展示了从原始语言模型到可用聊天机器人的渐进式改造过程：

\textbf{第一步：Role Prompting（角色提示）。} 在用户输入前添加 ``You are a helpful chatbot''，将模型引导到训练数据中``有帮助''的文本区域。效果略有改善，但模型同时扮演了对话双方。

\textbf{第二步：格式提示。} 使用类似电影剧本的格式 ``User: ...\textbackslash n Chatbot: ...'' 来提示模型区分对话角色。模型立即捕捉到了格式暗示，甚至给自己起了名字（``HelpBot''）。

\begin{knowledgebox}{为什么格式提示有效？}
模型的训练数据中大量存在格式化的对话数据（论坛帖子、电影剧本、客服记录等），这些数据通常以 ``User:'' / ``Assistant:'' 之类的标签区分说话人。通过使用相似的格式，我们把模型引导到了训练数据中对话格式文本所对应的概率空间区域。模型并不``理解''角色的概念，它只是在生成最可能跟在这种格式后面的文本。
\end{knowledgebox}

\textbf{第三步：截断输出。} 模型仍会同时生成双方的对话。解决方法很简单：只保留 ``Chatbot:'' 后面到下一个 ``User:'' 标签之间的文本，丢弃其余部分。

\textbf{第四步：构建对话 harness。} 编写一个对话管理程序，维护完整对话历史，每次将历史拼接到 prompt 中发送给模型，获取回复后追加到历史，循环往复。

\begin{warningbox}{模型并非在``对话''}
讲师特别强调，这整个过程``不是 Terminator 2: Rise of the Machines 的前奏''。模型同时生成双方对话不是因为它有自我意识，而是因为它的训练数据包含大量完整的对话文本（电影剧本等），它只是在做 next-word prediction。我们需要工程手段（格式提示 + 截断 + harness）来约束它的行为。
\end{warningbox}

\subsection{本章小结}

从原始语言模型到聊天机器人的关键步骤是：(1) 通过 role prompting 引导模型行为；(2) 通过格式提示帮助模型区分角色；(3) 通过截断和对话 harness 实现交互式多轮对话。整个过程的本质是利用 prompt 将模型引导到训练数据概率空间中我们期望的区域。

